\documentclass[../../main.tex]{subfiles}

% Directed Reading Program notes for October 1, 2026.

\begin{document}

\section{Tropical Arithmetic and Polynomials}

\subsection{Tropical Semiring}

\begin{defin}
    The \textit{min-plus tropical semiring} is
    \[\{\rr \cup \{\infty\},\oplus,\odot\},\]
    with tropical addition $a \oplus b = \min\{a,b\}$ and 
    tropical multiplication $a \odot b = a+b$.
\end{defin}

By the definition of a semiring, the following holds: 
\begin{prop}
    For any $x,y,z \in \rr\cup \{\infty\}$,
    \begin{itemize}
        \item (Commutativity)
        \[x\oplus y=y\oplus x \quad \text{and}\quad x\odot y=y\odot x.\]
        \item (Distributive law)
        \[x\odot(y\oplus z)=x\odot y\oplus x\odot z.\]
        \item (Identity elements)
        \[x\oplus \infty=x\quad \text{and}\quad x\odot 0=0.\]
    \end{itemize}

    Note that the existence of an additive inverse is not guaranteed. For example, $6\oplus x=7$ has no solution for $x$.
\end{prop}

\subsection{Tropical Polynomials}

\begin{defin}
    A \textit{tropical polynomial} is a finite linear combination of tropical monomials 
    \[f(x_1,\dots,x_n)=a\odot x_1^{i_1}x_2^{i_2}\cdots x_n^{i_n}\oplus b\odot x_1^{j_1}x_2^{j_2}\cdots x_n\oplus \cdots\]
    where the coefficients $a,b,\dots \in \rr$ and the exponents $i_1,j_1,\dots \in \zz_{\ge 0}$. 

    A tropical polynomial $f$ can be regarded as a function $f:\rr^n \to \rr$ by
    \[f(x_1,\dots,x_n)=\min(a+i_1x_1+\dots+i_nx_n,b+j_1x_1+\dots+j_nx_n,\dots)\]

    To use vectors $\xx=(x_1,\dots,x_n) \in \rr^n$ and $\ii=(i_1,\dots,i_n)\in \zz_{\ge 0}^n$,
    \[f(\xx)=\bigoplus_{\ii \in S}a_\ii \odot \xx^\ii,\]
    where $A\subseteq \zz_{\ge 0}^n$ is finite and 
    \[\xx^\ii=x_1^{i_1}\odot x_2^{i_2}\odot \cdots \odot x_n^{i_n}.\]

    In the regular arithmetic,
    \[f(\xx)=\min_{\ii \in S}\{a_\ii+\langle \ii,\xx\rangle\}.\]
\end{defin}

\begin{lem}
    Every tropical polynomial function $f:\rr^n \to \rr$ is
    \begin{itemize}
        \item continuous;
        \item piecewise-linear, with finitely many domains of linearity;
        \item concave.
    \end{itemize}
\end{lem}

\begin{proof}
    Continuity follows because a finite minimum of continuous functions is continuous.

    Since each monomial represents a linear function, taking the minimum of them gives a piecewise-linear function with no more domains of linearity than monomials.

    Write $f(\xx)=\min_{\ii \in A}\{a_\ii + \langle \ii,\xx\rangle\}$. For any $\xx,\yy \in \rr^n$ and $t \in [0,1]$,
    \begin{align*}
        f(t\xx+(1-t)\yy)
        &=\min_{\ii \in S}(a_\ii+t\langle \ii,\xx\rangle+(1-t)\langle \ii,\yy\rangle)\\
        &=\min_{\ii \in S}(t(a_\ii+\langle \ii,\xx) +(1-t)(a_\ii+\langle \ii,\yy\rangle)\\
        &\ge tf(\xx)+(1-t)f(\yy).
    \end{align*}

    Hence, $f$ is concave.
\end{proof}

\begin{lem} \label{lem:polynomial-pl graphs}
    The tropical polynomials in $n$ variables $x_1,\dots,x_n$ are precisely the piecewise-linear concave functions on $\rr^n$ with integer coefficients.
\end{lem}

\subsection{Fundamental Theorem of Tropical Algebra}

Every tropical polynomial function can be written uniquely as a tropical product of tropical linear functions. For example, 
\[x^2\oplus 17\odot x\oplus 2 \quad = \quad x^2\oplus 1\odot x\oplus 2 \quad =\quad (x\oplus 1)^2.\]

A factorization of a tropical polynomial gives the following properties.

\begin{prop}
    Let $f:\rr \to \rr$ be a tropical polynomial with one variable with the factorization
    \[f(x)=(x\oplus r_1)^{a_1}\odot (x\oplus r_1)^{a_2}\odot\cdots \odot (x\oplus r_n)^{a_n}\]
    where $r_1<r_2<\dots<r_n$.

    Then,
    \[f(x)=\begin{cases}
        \left(\sum_{j=1}^n a_j\right)x, & x\le r_1,\\
        \left(\sum_{j=i+1}^n a_j\right)x+\sum_{j=1}^i (a_jr_j), & r_i<x\le r_{i+1},\\
        \sum_{j=1}^n (a_jr_j), & x>r_n.
    \end{cases}\]

    Hence, the gradient at $x\neq r_i$ is
    \[n-\#\{i:r_i\le x\}.\]
\end{prop}

\begin{thm}[Fundamental Theorem of Tropical Algebra]
    Let $f:\rr \to \rr$ be given by 
    \[f(x)=\bigoplus_{i=0}^n a_ix^i=\min_{i=0}^n\{a_i+ix\}.\]
    Then, there exist $r_1,\dots,r_n$ such that
    \[p(x)=a_n\odot(x\oplus r_1)\odot \cdots \odot (x\oplus r_n)\]
    satisfying $f=p$ as functions $\rr\to \rr$.
\end{thm}

\begin{proof}[Proof sketch]
    We establish two bijections:
    \begin{enumerate}[label=(\roman*)]
        \item a bijection between such $p$'s and graphs ending with zero gradient when $x \to \infty$, gradient of $n$ when $x \to -\infty$, and non-increasing integer gradients;

        \item a bijection between such graphs and $f$'s.
    \end{enumerate}
\end{proof}

\subsection{Newton Polytopes}

\begin{defin}
    For any tropical polynomial
    \[f(\xx)=\bigoplus_{\ii \in S}a_\ii \odot \xx^\ii,\]
    the \textit{Newton polytope} of $f$ is
    \[\Delta_f=\conv (S)\subseteq \rr^n.\]
    
    It is the convex hull of the exponent vectors of the monomials in $f$.
\end{defin}

\begin{prop}
    For any two tropical polynomials $f$ and $g$,
    \[\Delta_{f\odot g}=\Delta_f+\Delta_g\]
    where $+$ denotes the Minkowski sum.
\end{prop}

\begin{example}
    Consider $f(x,y)=(x\oplus 0)(y\oplus 0)(xy\oplus 0)$.

    \begin{align*}
        f(x,y) &= x^2y^2 \oplus x^2y \oplus xy^2\oplus x\oplus y\oplus 0\\
        &= (xy\oplus x\oplus 0)(xy\oplus y\oplus 0).
    \end{align*}

    Its Newton polygon is as follows, which equals
    \[\conv((1,1),(1,0),(0,0))+\conv((1,1),(0,1),(0,0))\]
    or
    \[\conv((1,0),(0,0))+\conv((0,1),(0,0))+\conv((1,1),(0,0)).\]
    
    Also note that unique factorization breaks in this case, where $f(x,y)$ has two distinct irreducible factorization.
\end{example}

\end{document}
